% Part: proof-theory
% Chapter: cut-elimination
% Section: interpolation

\documentclass[../../../include/open-logic-section]{subfiles}

\begin{document}

\olfileid{pt}{cut}{itp}

\olsection{في لمّة ماهارا ومبرهنة كريغ في الاستيفاء}

مضمون مبرهنة الاستيفاء، وهي لوليم كريغ، أنه متى كان
$!A \Entails !B$ وجدت~!!a{sentence}~$!C$، تسمى
\emph{صيغة الاستيفاء}، تحقق~$!A \Entails !C$ و$!C \Entails !B$.
ويشترط في اختيار صيغة الاستيفاء~$!C$ أن تكون كل
!!{constant}s و!!{predicate}s الواقعة فيها واقعة في كل من~$!A$ و$!B$.
(ونفرض أنه لا~!!{function}s.) إذ لولا هذا الشرط لأغنت~$!A$
و$!B$ أنفسهما، على وجه لا يفيد، صيغتي استيفاء.

وتصح المبرهنة في منطق الرتبة الأولى الكلاسيكي، حتى وصفت بأنها «آخر خاصية
مهمة لمنطق الرتبة الأولى» اكتشفت. ومن منافعها تطبيقات في نظرية النماذج
والاستدلال الآلي. وأصل الداعي إليها وأول تطبيقها من فلسفة العلم؛ إذ يبحث
بها في الأوليات التي تقبل معها نظرية أو لا تقبل. ومنها تلزم مبرهنة بيث
في قابلية التعريف، ومؤداها أن النظرية إذا عرفت مفهومًا ضمنًا أمكنها أن
تعرفه تصريحًا.

ولمبرهنة الاستيفاء، كسائر نتائج كثيرة في المنطق الرياضي، براهين من جهة
نظرية النماذج وبراهين من جهة نظرية البرهان. وتمتاز الثانية بأنها لا تثبت
وجود صيغة استيفاء فحسب، بل تعطي خوارزمية تستخرجها من~!!a{proof} صوري
للمتتالية~$!A \Sequent !B$، مثل~!!a{proof} لها في
\Log{G3c}.

نرمز بـ~$\Lang L(!A)$ إلى مجموعة~!!{constant}s و!!{predicate}s الواقعة
في~$!A$، ونضع~$\Lang L({\OLId{greek-variable}{Gamma}}) = \bigcup \Setabs{\Lang L(!A)}{!A \in
{\OLId{greek-variable}{Gamma}}}$. والفرض الجاري في هذا القسم كله أن~!!{formula}s التي نعالجها
خالية من~!!{function}s.

\begin{thm}[مبرهنة كريغ للاستيفاء]\ollabel{thm:interpolation}
لتكن~!!{language}~$!A$ هي~$\Lang L(!A)$، ولغة~$!B$ هي
$\Lang L(!B)$. إذا كان~$\Log{G3c} \Proves !A \Sequent !B$، وُجدت
!!a{formula}~$!C$ في~!!{language}~$\Lang L(!C) \subseteq \Lang
L(!A) \cap \Lang L(!B)$ بحيث~$\Log{G3c} \Proves !A \Sequent !C$
و$\Log{G3c} \Proves !C \Sequent !B$.
\end{thm}

نثبت أولًا صورة أعم من مبرهنة الاستيفاء، لننتفع ببنية~!!{proof}s في
\Log{G3c}. فنقسم مقدّم كل متتالية وتاليها قسمين، ونكتب~$\maeh{{\OLId{greek-variable}{Gamma}}_{\OLMathNumeral{1}}}{{\OLId{greek-variable}{Gamma}}_{\OLMathNumeral{2}}} \Sequent
\maeh{{\OLId{greek-variable}{Delta}}_{\OLMathNumeral{1}}}{{\OLId{greek-variable}{Delta}}_{\OLMathNumeral{2}}}$ للدلالة على أن~${\OLId{greek-variable}{Gamma}}_{\OLMathNumeral{1}}$ و${\OLId{greek-variable}{Delta}}_{\OLMathNumeral{1}}$
ينتميان إلى الجزء الأول، وأن~${\OLId{greek-variable}{Gamma}}_{\OLMathNumeral{2}}$ و${\OLId{greek-variable}{Delta}}_{\OLMathNumeral{2}}$ ينتميان إلى الجزء
الثاني. ومن اللمّة الآتية تنتج مبرهنة الاستيفاء من فورها.

\begin{lem}[لمّة ماهارا]\ollabel{lem:maehara} إذا كان~$\Log{G3c}
\Proves \maeh{{\OLId{greek-variable}{Gamma}}_{\OLMathNumeral{1}}}{{\OLId{greek-variable}{Gamma}}_{\OLMathNumeral{2}}} \Sequent \maeh{{\OLId{greek-variable}{Delta}}_{\OLMathNumeral{1}}}{{\OLId{greek-variable}{Delta}}_{\OLMathNumeral{2}}}$،
وُجدت~!!a{formula}~$!C$ في~!!{language}~$\Lang L(!C)
\subseteq \Lang L({\OLId{greek-variable}{Gamma}}_{\OLMathNumeral{1}},
{\OLId{greek-variable}{Delta}}_{\OLMathNumeral{1}}) \cap \Lang L({\OLId{greek-variable}{Gamma}}_{\OLMathNumeral{2}}, {\OLId{greek-variable}{Delta}}_{\OLMathNumeral{2}})$ بحيث~$\Log{G3c}
\Proves {\OLId{greek-variable}{Gamma}}_{\OLMathNumeral{1}} \Sequent {\OLId{greek-variable}{Delta}}_{\OLMathNumeral{1}}, !C$ و$\Log{G3c} \Proves !C,
{\OLId{greek-variable}{Gamma}}_{\OLMathNumeral{2}} \Sequent {\OLId{greek-variable}{Delta}}_{\OLMathNumeral{2}}$.
\end{lem}

\begin{proof}
  لنفرض أن~$\pi \Proves \maeh{{\OLId{greek-variable}{Gamma}}_{\OLMathNumeral{1}}}{{\OLId{greek-variable}{Gamma}}_{\OLMathNumeral{2}}} \Sequent
  \maeh{{\OLId{greek-variable}{Delta}}_{\OLMathNumeral{1}}}{{\OLId{greek-variable}{Delta}}_{\OLMathNumeral{2}}}$. نبرهن المطلوب بالاستقراء على
  ${\OLId{latin-ordinary}{n}}=\pheight{\pi}$.

  إذا كان~${\OLId{latin-ordinary}{n}}={\OLMathNumeral{0}}$، كانت~$\maeh{{\OLId{greek-variable}{Gamma}}_{\OLMathNumeral{1}}}{{\OLId{greek-variable}{Gamma}}_{\OLMathNumeral{2}}} \Sequent
  \maeh{{\OLId{greek-variable}{Delta}}_{\OLMathNumeral{1}}}{{\OLId{greek-variable}{Delta}}_{\OLMathNumeral{2}}}$ مسلّمة، وظهرت~!!{formula} ذرية ما~$!A$
  في كل من~${\OLId{greek-variable}{Gamma}}_{\OLMathNumeral{1}}, {\OLId{greek-variable}{Gamma}}_{\OLMathNumeral{2}}$ و${\OLId{greek-variable}{Delta}}_{\OLMathNumeral{1}}, {\OLId{greek-variable}{Delta}}_{\OLMathNumeral{2}}$. ولدينا أربع
  حالات بحسب انتماء~$!A$ إلى الجزء الأول أو الثاني في اليسار، وانتمائها
  إلى الجزء الأول أو الثاني في اليمين.
  \begin{enumerate}
    \item $!A \in {\OLId{greek-variable}{Gamma}}_{\OLMathNumeral{1}}$ و$!A \in {\OLId{greek-variable}{Delta}}_{\OLMathNumeral{1}}$. نريد~$!C$ بحيث
    تكون كل من~${\OLId{greek-variable}{Gamma}}_{\OLMathNumeral{1}} \Sequent {\OLId{greek-variable}{Delta}}_{\OLMathNumeral{1}}, !C$ و$!C, {\OLId{greek-variable}{Gamma}}_{\OLMathNumeral{2}}
    \Sequent {\OLId{greek-variable}{Delta}}_{\OLMathNumeral{2}}$~!!{provable}. الأولى مسلّمة بالفعل أيًا كان
    اختيار~$!C$، لأن~$!A$ تظهر في اليسار واليمين. والطريقة الوحيدة
    لضمان كون~$!C, {\OLId{greek-variable}{Gamma}}_{\OLMathNumeral{2}} \Sequent {\OLId{greek-variable}{Delta}}_{\OLMathNumeral{2}}$~!!{provable} هي أن
    نضع~$!C \ident \lfalse$.
    \item $!A \in {\OLId{greek-variable}{Gamma}}_{\OLMathNumeral{1}}$ و$!A \in {\OLId{greek-variable}{Delta}}_{\OLMathNumeral{2}}$. عندئذ تكون
    ${\OLId{greek-variable}{Gamma}}_{\OLMathNumeral{1}} \Sequent {\OLId{greek-variable}{Delta}}_{\OLMathNumeral{1}}, !A$ و$!A, {\OLId{greek-variable}{Gamma}}_{\OLMathNumeral{2}} \Sequent
    {\OLId{greek-variable}{Delta}}_{\OLMathNumeral{2}}$ مسلّمتين، فيمكن اختيار~$!C \ident !A$.
    \item $!A \in {\OLId{greek-variable}{Gamma}}_{\OLMathNumeral{2}}$ و$!A \in {\OLId{greek-variable}{Delta}}_{\OLMathNumeral{1}}$. عندئذ تكون
    $!A, {\OLId{greek-variable}{Gamma}}_{\OLMathNumeral{1}} \Sequent {\OLId{greek-variable}{Delta}}_{\OLMathNumeral{1}}$ و${\OLId{greek-variable}{Gamma}}_{\OLMathNumeral{2}} \Sequent {\OLId{greek-variable}{Delta}}_{\OLMathNumeral{2}},
    !A$ مسلّمتين، غير أن~$!A$ تقع فيهما في الجهتين الخطأ، فلا تكون هي صيغة
    الاستيفاء. غير أننا نستطيع باستعمال~$\RightR{\lnot}$ و
    $\LeftR{\lnot}$ أن~!!{prove} المتتاليتين~${\OLId{greek-variable}{Gamma}}_{\OLMathNumeral{1}}
    \Sequent {\OLId{greek-variable}{Delta}}_{\OLMathNumeral{1}}, \lnot!A$ و$\lnot !A, {\OLId{greek-variable}{Gamma}}_{\OLMathNumeral{2}} \Sequent
    {\OLId{greek-variable}{Delta}}_{\OLMathNumeral{2}}$.
    \item $!A \in {\OLId{greek-variable}{Gamma}}_{\OLMathNumeral{2}}$ و$!A \in {\OLId{greek-variable}{Delta}}_{\OLMathNumeral{2}}$. تمرين.
  \end{enumerate}
  ونترك تمرينًا حالتي كون~$\maeh{{\OLId{greek-variable}{Gamma}}_{\OLMathNumeral{1}}}{{\OLId{greek-variable}{Gamma}}_{\OLMathNumeral{2}}} \Sequent
  \maeh{{\OLId{greek-variable}{Delta}}_{\OLMathNumeral{1}}}{{\OLId{greek-variable}{Delta}}_{\OLMathNumeral{2}}}$ مسلّمة بسبب~$\lfalse \in {\OLId{greek-variable}{Gamma}}_{\OLMathNumeral{1}}$ أو
  $\lfalse \in {\OLId{greek-variable}{Gamma}}_{\OLMathNumeral{2}}$.

  إذا كان~${\OLId{latin-ordinary}{n}}>{\OLMathNumeral{0}}$، انتهى~$\pi$ باستدلال منطقي. وعلينا تمييز الحالات بحسب
  القاعدة المستخدمة والجزء الذي تنتمي إليه الصيغة الرئيسة. وفي كل حالة
  نستطيع استعمال فرض الاستقراء القائل إن اللمّة تصدق لكل~!!{proof}s
  $\pi_{\OLMathNumeral{1}}$ بحيث~$\pheight{\pi_{\OLMathNumeral{1}}} < {\OLId{latin-ordinary}{n}}$، وبأي تقسيم لمتتالية نهاية
  $\pi_{\OLMathNumeral{1}}$ إلى جزأين. فلنتأمل بعض الحالات المحددة.
  \begin{enumerate}
    \item ينتهي~$\pi$ بـ$\RightR{\lnot}$، وصيغته الرئيسة~$\lnot
    !A$. ولدينا حالتان: بحسب انتماء~$\lnot !A$ إلى الجزء الأول أو
    الثاني، ينتهي~$\pi$ بأحد الشكلين
    \[
    \AxiomC{}
    \RightLabel{$\pi_{\OLMathNumeral{1}}$}
    \Deduce$\maeh{!A, {\OLId{greek-variable}{Gamma}}_{\OLMathNumeral{1}}}{{\OLId{greek-variable}{Gamma}}_{\OLMathNumeral{2}}} \fCenter
      \maeh{{\OLId{greek-variable}{Delta}}_{\OLMathNumeral{1}}}{{\OLId{greek-variable}{Delta}}_{\OLMathNumeral{2}}}$
    \RightLabel{\RightR{\lnot}}
    \UnaryInf$\maeh{{\OLId{greek-variable}{Gamma}}_{\OLMathNumeral{1}}}{{\OLId{greek-variable}{Gamma}}_{\OLMathNumeral{2}}} \fCenter
      \maeh{{\OLId{greek-variable}{Delta}}_{\OLMathNumeral{1}}, \lnot !A}{{\OLId{greek-variable}{Delta}}_{\OLMathNumeral{2}}}$
    \DisplayProof
    \]
    أو
    \[
    \AxiomC{}
    \RightLabel{$\pi_{\OLMathNumeral{1}}$}
    \Deduce$\maeh{{\OLId{greek-variable}{Gamma}}_{\OLMathNumeral{1}}}{!A, {\OLId{greek-variable}{Gamma}}_{\OLMathNumeral{2}}} \fCenter
      \maeh{{\OLId{greek-variable}{Delta}}_{\OLMathNumeral{1}}}{{\OLId{greek-variable}{Delta}}_{\OLMathNumeral{2}}}$
    \RightLabel{\RightR{\lnot}}
    \UnaryInf$\maeh{{\OLId{greek-variable}{Gamma}}_{\OLMathNumeral{1}}}{{\OLId{greek-variable}{Gamma}}_{\OLMathNumeral{2}}} \fCenter
      \maeh{{\OLId{greek-variable}{Delta}}_{\OLMathNumeral{1}}}{{\OLId{greek-variable}{Delta}}_{\OLMathNumeral{2}}, \lnot !A}$
    \DisplayProof
    \]
    لاحظ أننا إذا وضعنا~!!{formula} الرئيسة~$\lnot !A$ في الجزء الأول،
    وضعنا~!!{formula} الفعالة~$!A$ في الجزء الأول من المقدمة كذلك. وهذا
    اختيار متاح لنا. وكان فرض الاستقراء سيصدق كذلك لو وضعنا الصيغة
    الفعالة في الجزء الآخر، لكننا عندئذ لن نستطيع إيجاد صيغة استيفاء
    (جرّب ذلك تمرينًا).

    تأمل أولًا الوضع الأول، حيث تنتمي~$\lnot !A$ إلى الجزء الأول. بفرض
    الاستقراء توجد صيغة استيفاء لمتتالية نهاية~$\pi_{\OLMathNumeral{1}}$، أي توجد
    !!a{formula}~$!C_{\OLMathNumeral{1}}$ بحيث تكون كل من
    \begin{align*}
    !A, {\OLId{greek-variable}{Gamma}}_{\OLMathNumeral{1}} & \Sequent {\OLId{greek-variable}{Delta}}_{\OLMathNumeral{1}}, !C_{\OLMathNumeral{1}} &&\text{و} &
    !C_{\OLMathNumeral{1}}, {\OLId{greek-variable}{Gamma}}_{\OLMathNumeral{2}} & \Sequent {\OLId{greek-variable}{Delta}}_{\OLMathNumeral{2}}
    \intertext{قابلة للبرهنة. وما نحتاج إليه~!!a{formula}~$!C$ بحيث
    تكون المتتاليتان الآتيتان~!!{provable}:}
    {\OLId{greek-variable}{Gamma}}_{\OLMathNumeral{1}} & \Sequent {\OLId{greek-variable}{Delta}}_{\OLMathNumeral{1}}, \lnot !A, !C &&\text{و} &
    !C, {\OLId{greek-variable}{Gamma}}_{\OLMathNumeral{2}} & \Sequent {\OLId{greek-variable}{Delta}}_{\OLMathNumeral{2}}
    \end{align*}
    ومن الواضح أنه يكفينا وضع~$!C \ident !C_{\OLMathNumeral{1}}$: فقد أخبرنا فرض
    الاستقراء أن المتتالية اليمنى~!!{provable}، وتنتج المتتالية اليسرى
    من المتتالية التي أعطاها فرض الاستقراء بتطبيق~$\RightR{\lnot}$.
    ويعطينا فرض الاستقراء كذلك~$\Lang L(!C_{\OLMathNumeral{1}}) \subseteq \Lang L(!A,
    {\OLId{greek-variable}{Gamma}}_{\OLMathNumeral{1}}, {\OLId{greek-variable}{Delta}}_{\OLMathNumeral{1}}) \cap \Lang L({\OLId{greek-variable}{Gamma}}_{\OLMathNumeral{2}}, {\OLId{greek-variable}{Delta}}_{\OLMathNumeral{2}})$. وبما أن
    $\Lang L(\lnot !A) = \Lang L(!A)$ و$\Lang L(!C) = \Lang L(!C_{\OLMathNumeral{1}})$،
    يكون~$\Lang L(!C_{\OLMathNumeral{1}}) \subseteq \Lang L({\OLId{greek-variable}{Gamma}}_{\OLMathNumeral{1}}, {\OLId{greek-variable}{Delta}}_{\OLMathNumeral{1}},
    \lnot !A) \cap \Lang L({\OLId{greek-variable}{Gamma}}_{\OLMathNumeral{2}}, {\OLId{greek-variable}{Delta}}_{\OLMathNumeral{2}})$.

    أما في الوضع الثاني فالحال مختلفة قليلًا، إذ تنتمي~$\lnot !A$ إلى
    الجزء الثاني. نضع~!!{formula} الفعالة في المقدمة في الجزء الثاني
    أيضًا، ونطبق فرض الاستقراء فنحصل على
    \begin{align*}
    {\OLId{greek-variable}{Gamma}}_{\OLMathNumeral{1}} & \Sequent {\OLId{greek-variable}{Delta}}_{\OLMathNumeral{1}}, !C_{\OLMathNumeral{1}} &&\text{و} &
    !C_{\OLMathNumeral{1}}, !A, {\OLId{greek-variable}{Gamma}}_{\OLMathNumeral{2}} & \Sequent {\OLId{greek-variable}{Delta}}_{\OLMathNumeral{2}}
    \intertext{وهما قابلتان للبرهنة. وإذا طبقنا قاعدة~$\RightR{\lnot}$
    مرة أخرى على الثانية، حصلنا على~!!{proof}s لـ}
    {\OLId{greek-variable}{Gamma}}_{\OLMathNumeral{1}} & \Sequent {\OLId{greek-variable}{Delta}}_{\OLMathNumeral{1}}, !C_{\OLMathNumeral{1}} &&\text{و} &
    !C_{\OLMathNumeral{1}}, {\OLId{greek-variable}{Gamma}}_{\OLMathNumeral{2}} & \Sequent {\OLId{greek-variable}{Delta}}_{\OLMathNumeral{2}}, \lnot !A
    \end{align*}
    وهاتان بالضبط المتتاليتان اللتان نريد أن تكونا~!!{provable} إذا كانت
    $!C_{\OLMathNumeral{1}}$ صيغة استيفاء؛ فنضع مرة أخرى~$!C \ident !C_{\OLMathNumeral{1}}$. وكما سبق يكون
    $\Lang L(!C_{\OLMathNumeral{1}}) \subseteq \Lang L({\OLId{greek-variable}{Gamma}}_{\OLMathNumeral{1}}, {\OLId{greek-variable}{Delta}}_{\OLMathNumeral{1}}) \cap
    \Lang L({\OLId{greek-variable}{Gamma}}_{\OLMathNumeral{2}}, {\OLId{greek-variable}{Delta}}_{\OLMathNumeral{2}}, \lnot !A)$.

    \item ينتهي~$\pi$ بـ$\RightR{\lif}$، وصيغته الرئيسة
    !!{formula}~$!A
    \lif !B$. ولدينا حالتان أيضًا، بحسب انتماء~$!A \lif !B$ إلى الجزء
    الأول أو الثاني. ينتهي~!!{proof}~$\pi$ بأحد الشكلين
    \[
    \AxiomC{}
    \RightLabel{$\pi_{\OLMathNumeral{1}}$}
    \Deduce$\maeh{!A, {\OLId{greek-variable}{Gamma}}_{\OLMathNumeral{1}}}{{\OLId{greek-variable}{Gamma}}_{\OLMathNumeral{2}}} \fCenter
      \maeh{{\OLId{greek-variable}{Delta}}_{\OLMathNumeral{1}}, !B}{{\OLId{greek-variable}{Delta}}_{\OLMathNumeral{2}}}$
    \RightLabel{\RightR{\lif}}
    \UnaryInf$\maeh{{\OLId{greek-variable}{Gamma}}_{\OLMathNumeral{1}}}{{\OLId{greek-variable}{Gamma}}_{\OLMathNumeral{2}}} \fCenter
      \maeh{{\OLId{greek-variable}{Delta}}_{\OLMathNumeral{1}}, !A \lif !B}{{\OLId{greek-variable}{Delta}}_{\OLMathNumeral{2}}}$
    \DisplayProof
    \]
    أو
    \[
    \AxiomC{}
    \RightLabel{$\pi_{\OLMathNumeral{1}}$}
    \Deduce$\maeh{{\OLId{greek-variable}{Gamma}}_{\OLMathNumeral{1}}}{!A, {\OLId{greek-variable}{Gamma}}_{\OLMathNumeral{2}}} \fCenter
      \maeh{{\OLId{greek-variable}{Delta}}_{\OLMathNumeral{1}}}{{\OLId{greek-variable}{Delta}}_{\OLMathNumeral{2}}, !B}$
    \RightLabel{\RightR{\lif}}
    \UnaryInf$\maeh{{\OLId{greek-variable}{Gamma}}_{\OLMathNumeral{1}}}{{\OLId{greek-variable}{Gamma}}_{\OLMathNumeral{2}}} \fCenter
      \maeh{{\OLId{greek-variable}{Delta}}_{\OLMathNumeral{1}}}{{\OLId{greek-variable}{Delta}}_{\OLMathNumeral{2}}, !A \lif !B}$
    \DisplayProof
    \]
    وقد وضعنا مرة أخرى~!!{formula}s الفعالة في الجزء نفسه من المقدمة
    الذي تنتمي إليه~!!{formula} الرئيسة في النتيجة. وفي الوضع الأول
    نطبق فرض الاستقراء فنحصل على~!!{proof}s لـ
    \begin{align*}
    !A, {\OLId{greek-variable}{Gamma}}_{\OLMathNumeral{1}} & \Sequent {\OLId{greek-variable}{Delta}}_{\OLMathNumeral{1}}, !B, !C_{\OLMathNumeral{1}} &&\text{و} & !C_{\OLMathNumeral{1}}, {\OLId{greek-variable}{Gamma}}_{\OLMathNumeral{2}} & \Sequent {\OLId{greek-variable}{Delta}}_{\OLMathNumeral{2}}
    \intertext{والمتتالية اليسرى مقدمة مناسبة لـ$\RightR{\lif}$؛ ومن
    ثم تكون المتتاليتان الآتيتان~!!{provable}:}
    {\OLId{greek-variable}{Gamma}}_{\OLMathNumeral{1}} & \Sequent {\OLId{greek-variable}{Delta}}_{\OLMathNumeral{1}}, !A \lif !B, !C_{\OLMathNumeral{1}} &&\text{و} & !C_{\OLMathNumeral{1}}, {\OLId{greek-variable}{Gamma}}_{\OLMathNumeral{2}} & \Sequent {\OLId{greek-variable}{Delta}}_{\OLMathNumeral{2}}
    \end{align*}
    ومؤدى ذلك أن~$!C_{\OLMathNumeral{1}}$ صيغة استيفاء لمتتالية نهاية~$\pi$ أيضًا،
    فيمكن أن نضع~$!C \ident !C_{\OLMathNumeral{1}}$ مرة أخرى.

    وتعالج الحالة الأخرى على هذا القياس.

    \item ينتهي~$\pi$ بـ$\LeftR{\lif}$، وصيغته الرئيسة
    !!{formula}~$!A
    \lif !B$. إذا انتمت~$!A \lif !B$ إلى الجزء الأول، انتهى
    !!{proof}~$\pi$ بالشكل
    \[
    \AxiomC{}
    \RightLabel{$\pi_{\OLMathNumeral{1}}$}
    \Deduce$\maeh{{\OLId{greek-variable}{Gamma}}_{\OLMathNumeral{1}}}{{\OLId{greek-variable}{Gamma}}_{\OLMathNumeral{2}}} \fCenter
      \maeh{{\OLId{greek-variable}{Delta}}_{\OLMathNumeral{1}}, !A}{{\OLId{greek-variable}{Delta}}_{\OLMathNumeral{2}}}$
    \AxiomC{}
    \RightLabel{$\pi_{\OLMathNumeral{2}}$}
    \Deduce$\maeh{!B, {\OLId{greek-variable}{Gamma}}_{\OLMathNumeral{1}}}{{\OLId{greek-variable}{Gamma}}_{\OLMathNumeral{2}}} \fCenter
      \maeh{{\OLId{greek-variable}{Delta}}_{\OLMathNumeral{1}}}{{\OLId{greek-variable}{Delta}}_{\OLMathNumeral{2}}}$
    \RightLabel{\LeftR{\lif}}
    \BinaryInf$\maeh{!A \lif !B, {\OLId{greek-variable}{Gamma}}_{\OLMathNumeral{1}}}{{\OLId{greek-variable}{Gamma}}_{\OLMathNumeral{2}}} \fCenter
      \maeh{{\OLId{greek-variable}{Delta}}_{\OLMathNumeral{1}}}{{\OLId{greek-variable}{Delta}}_{\OLMathNumeral{2}}}$
    \DisplayProof
    \]
    ويعطينا فرض الاستقراء الآن أربع متتاليات~!!{provable}: اثنتان لصيغة
    الاستيفاء~$!C_{\OLMathNumeral{1}}$ للمقدمة اليسرى، واثنتان لصيغة الاستيفاء~$!C_{\OLMathNumeral{2}}$
    للمقدمة الثانية:
    \begin{align*}
    {\OLId{greek-variable}{Gamma}}_{\OLMathNumeral{1}} & \Sequent {\OLId{greek-variable}{Delta}}_{\OLMathNumeral{1}}, !A, !C_{\OLMathNumeral{1}} &&\text{(1)} &
    !C_{\OLMathNumeral{1}}, {\OLId{greek-variable}{Gamma}}_{\OLMathNumeral{2}} & \Sequent {\OLId{greek-variable}{Delta}}_{\OLMathNumeral{2}} && \text{(2)}\\
    !B, {\OLId{greek-variable}{Gamma}}_{\OLMathNumeral{1}} & \Sequent {\OLId{greek-variable}{Delta}}_{\OLMathNumeral{1}}, !C_{\OLMathNumeral{2}} &&\text{(3)} &
    !C_{\OLMathNumeral{2}}, {\OLId{greek-variable}{Gamma}}_{\OLMathNumeral{2}} & \Sequent {\OLId{greek-variable}{Delta}}_{\OLMathNumeral{2}} && \text{(4)}
    \end{align*}
    ويمكن أن تكون المتتاليتان (\OLClassicalDigits{1}) و(\OLClassicalDigits{3}) مقدمتي~$\LeftR{\lif}$ إذا طبقنا
    الإضعاف، فأضفنا~$!C_{\OLMathNumeral{2}}$ إلى يمين (\OLClassicalDigits{1}) و$!C_{\OLMathNumeral{1}}$ إلى يمين (\OLClassicalDigits{3}):
    \[
    \AxiomC{}
    \Deduce${\OLId{greek-variable}{Gamma}}_{\OLMathNumeral{1}} \fCenter {\OLId{greek-variable}{Delta}}_{\OLMathNumeral{1}}, !A, !C_{\OLMathNumeral{1}}, !C_{\OLMathNumeral{2}}$
    \AxiomC{}
    \Deduce$!B, {\OLId{greek-variable}{Gamma}}_{\OLMathNumeral{1}} \fCenter {\OLId{greek-variable}{Delta}}_{\OLMathNumeral{1}}, !C_{\OLMathNumeral{1}}, !C_{\OLMathNumeral{2}}$
    \RightLabel{\LeftR{\lif}}
    \BinaryInf$!A \lif !B, {\OLId{greek-variable}{Gamma}}_{\OLMathNumeral{1}} \fCenter {\OLId{greek-variable}{Delta}}_{\OLMathNumeral{1}}, !C_{\OLMathNumeral{1}}, !C_{\OLMathNumeral{2}}$
    \DisplayProof
    \]
    وبتطبيق~$\RightR{\lor}$ ينتج لنا المتتالية
    \[
    !A \lif !B, {\OLId{greek-variable}{Gamma}}_{\OLMathNumeral{1}} \fCenter {\OLId{greek-variable}{Delta}}_{\OLMathNumeral{1}}, !C_{\OLMathNumeral{1}} \lor !C_{\OLMathNumeral{2}}.
    \]
    وهذا يرجح أن تكون~$!C_{\OLMathNumeral{1}} \lor !C_{\OLMathNumeral{2}}$ صيغة الاستيفاء عندئذ.
    وهي كذلك بالفعل، إذ نستطيع أن~!!{prove} المتتالية الأخرى المطلوبة
    انطلاقًا من المتتاليتين الباقيتين (\OLClassicalDigits{2}) و(\OLClassicalDigits{4}):
    \[
    \AxiomC{}
    \Deduce$!C_{\OLMathNumeral{1}}, {\OLId{greek-variable}{Gamma}}_{\OLMathNumeral{2}} \fCenter {\OLId{greek-variable}{Delta}}_{\OLMathNumeral{2}}$
    \AxiomC{}
    \Deduce$!C_{\OLMathNumeral{2}}, {\OLId{greek-variable}{Gamma}}_{\OLMathNumeral{2}} \fCenter {\OLId{greek-variable}{Delta}}_{\OLMathNumeral{2}}$
    \RightLabel{\LeftR{\lor}}
    \BinaryInf$!C_{\OLMathNumeral{1}} \lor !C_{\OLMathNumeral{2}}, {\OLId{greek-variable}{Gamma}}_{\OLMathNumeral{2}} \fCenter {\OLId{greek-variable}{Delta}}_{\OLMathNumeral{2}}$
    \DisplayProof
    \]
    يبقى التحقق من أن~$!C_{\OLMathNumeral{1}} \lor !C_{\OLMathNumeral{2}}$ تنتمي إلى اللغة الصحيحة. يعطينا
    فرض الاستقراء
    \begin{align*}
    \Lang L(!C_{\OLMathNumeral{1}}) & \subseteq
    \Lang L({\OLId{greek-variable}{Gamma}}_{\OLMathNumeral{1}}, {\OLId{greek-variable}{Delta}}_{\OLMathNumeral{1}}, !A) \cap \Lang L({\OLId{greek-variable}{Gamma}}_{\OLMathNumeral{2}}, {\OLId{greek-variable}{Delta}}_{\OLMathNumeral{2}}) &\text{و}\\
    \Lang L(!C_{\OLMathNumeral{2}}) & \subseteq
    \Lang L(!B, {\OLId{greek-variable}{Gamma}}_{\OLMathNumeral{1}}, {\OLId{greek-variable}{Delta}}_{\OLMathNumeral{1}}) \cap \Lang L({\OLId{greek-variable}{Gamma}}_{\OLMathNumeral{2}}, {\OLId{greek-variable}{Delta}}_{\OLMathNumeral{2}})
    \intertext{ولما كان}
    \Lang L(!A) & \subseteq \Lang L(!A \lif !B) &\text{و}\\
    \Lang L(!B) & \subseteq \Lang L(!A \lif !B)
    \intertext{كان كل من}
    \Lang L(!C_{\OLMathNumeral{1}}) & \subseteq
    \Lang L({\OLId{greek-variable}{Gamma}}_{\OLMathNumeral{1}}, {\OLId{greek-variable}{Delta}}_{\OLMathNumeral{1}}, !A \lif !B) \cap \Lang L({\OLId{greek-variable}{Gamma}}_{\OLMathNumeral{2}}, {\OLId{greek-variable}{Delta}}_{\OLMathNumeral{2}}) &\text{و}\\
    \Lang L(!C_{\OLMathNumeral{2}}) & \subseteq
    \Lang L({\OLId{greek-variable}{Gamma}}_{\OLMathNumeral{1}}, {\OLId{greek-variable}{Delta}}_{\OLMathNumeral{1}}, !A \lif !B) \cap \Lang L({\OLId{greek-variable}{Gamma}}_{\OLMathNumeral{2}}, {\OLId{greek-variable}{Delta}}_{\OLMathNumeral{2}}).
    \intertext{وبالتالي، لما كان~$\Lang L (!C_{\OLMathNumeral{1}} \lor !C_{\OLMathNumeral{2}}) = {}$}
    \Lang L(!C_{\OLMathNumeral{1}}) \cup \Lang L(!C_{\OLMathNumeral{2}}) & \subseteq
    \Lang L({\OLId{greek-variable}{Gamma}}_{\OLMathNumeral{1}}, {\OLId{greek-variable}{Delta}}_{\OLMathNumeral{1}}, !A \lif !B) \cap \Lang L({\OLId{greek-variable}{Gamma}}_{\OLMathNumeral{2}}, {\OLId{greek-variable}{Delta}}_{\OLMathNumeral{2}}),
    \end{align*}
    وهو المطلوب.

    وإذا انتمت~$!A \lif !B$ إلى الجزء الثاني اختلف الوضع، لكنه يعالج
    على هذا القياس. ويعطينا فرض الاستقراء مرة أخرى أربع متتاليات
    !!{provable}:
    \begin{align*}
    {\OLId{greek-variable}{Gamma}}_{\OLMathNumeral{1}} & \Sequent {\OLId{greek-variable}{Delta}}_{\OLMathNumeral{1}}, !C_{\OLMathNumeral{1}} &&\text{(1)} &
    !C_{\OLMathNumeral{1}}, {\OLId{greek-variable}{Gamma}}_{\OLMathNumeral{2}} & \Sequent {\OLId{greek-variable}{Delta}}_{\OLMathNumeral{2}}, !A && \text{(2)}\\
    {\OLId{greek-variable}{Gamma}}_{\OLMathNumeral{1}} & \Sequent {\OLId{greek-variable}{Delta}}_{\OLMathNumeral{1}}, !C_{\OLMathNumeral{2}} &&\text{(3)} &
    !C_{\OLMathNumeral{2}}, !B, {\OLId{greek-variable}{Gamma}}_{\OLMathNumeral{2}} & \Sequent {\OLId{greek-variable}{Delta}}_{\OLMathNumeral{2}} && \text{(4)}
    \end{align*}
    والآن تكون المتتاليتان (\OLClassicalDigits{2}) و(\OLClassicalDigits{4})، بعد إضافة بعض~!!{formula}s
    بالإضعاف، مقدمتي~$\LeftR{\lif}$:
    \[
    \AxiomC{}
    \Deduce$!C_{\OLMathNumeral{1}}, !C_{\OLMathNumeral{2}}, {\OLId{greek-variable}{Gamma}}_{\OLMathNumeral{2}} \fCenter {\OLId{greek-variable}{Delta}}_{\OLMathNumeral{2}}, !A$
    \AxiomC{}
    \Deduce$!B, !C_{\OLMathNumeral{1}}, !C_{\OLMathNumeral{2}}, {\OLId{greek-variable}{Gamma}}_{\OLMathNumeral{2}} \fCenter {\OLId{greek-variable}{Delta}}_{\OLMathNumeral{2}}$
    \RightLabel{\LeftR{\lif}}
    \BinaryInf$!A \lif !B, !C_{\OLMathNumeral{1}}, !C_{\OLMathNumeral{2}}, {\OLId{greek-variable}{Gamma}}_{\OLMathNumeral{2}} \fCenter {\OLId{greek-variable}{Delta}}_{\OLMathNumeral{2}}$
    \DisplayProof
    \]
    ونريد مرة أخرى تحويل هذه إلى متتالية فيها~!!{formula} واحدة مبنية من
    $!C_{\OLMathNumeral{1}}$ و$!C_{\OLMathNumeral{2}}$، لكننا نحتاج إليها هذه المرة في المقدّم لأننا نتعامل
    الآن مع الجزء الثاني. ويؤدي تطبيق~$\LeftR{\land}$ الغرض؛ فينبغي أن
    نختار~$!C \ident (!C_{\OLMathNumeral{1}} \land !C_{\OLMathNumeral{2}})$ صيغة للاستيفاء. ومن حسن الحظ
    أننا نستطيع استعمال المتتاليتين الباقيتين (\OLClassicalDigits{1}) و(\OLClassicalDigits{3}) كي~!!{prove}
    المتتالية المقابلة للجزء الآخر:
    \[
    \AxiomC{}
    \Deduce${\OLId{greek-variable}{Gamma}}_{\OLMathNumeral{1}} \fCenter {\OLId{greek-variable}{Delta}}_{\OLMathNumeral{1}}, !C_{\OLMathNumeral{1}}$
    \AxiomC{}
    \Deduce${\OLId{greek-variable}{Gamma}}_{\OLMathNumeral{1}} \fCenter {\OLId{greek-variable}{Delta}}_{\OLMathNumeral{1}}, !C_{\OLMathNumeral{2}}$
    \RightLabel{\RightR{\land}}
    \BinaryInf${\OLId{greek-variable}{Gamma}}_{\OLMathNumeral{1}} \fCenter {\OLId{greek-variable}{Delta}}_{\OLMathNumeral{1}}, !C_{\OLMathNumeral{1}} \land !C_{\OLMathNumeral{2}}$
    \DisplayProof
    \]
    وكما سبق، يكون~$\Lang L(!C_{\OLMathNumeral{1}} \land !C_{\OLMathNumeral{2}}) \subseteq \Lang
    L({\OLId{greek-variable}{Gamma}}_{\OLMathNumeral{1}}, {\OLId{greek-variable}{Delta}}_{\OLMathNumeral{1}}) \cap \Lang L({\OLId{greek-variable}{Gamma}}_{\OLMathNumeral{2}}, {\OLId{greek-variable}{Delta}}_{\OLMathNumeral{2}}, !A \lif
    !B)$.

    \item ينتهي~$\pi$ بـ$\RightR{\lforall}$، وصيغته الرئيسة
    !!{formula}~$\lforall[{\OLId{latin-ordinary}{x}}][!A({\OLId{latin-ordinary}{x}})]$:
    \[
    \AxiomC{}
    \RightLabel{$\pi_{\OLMathNumeral{1}}$}
    \Deduce$\maeh{{\OLId{greek-variable}{Gamma}}_{\OLMathNumeral{1}}}{{\OLId{greek-variable}{Gamma}}_{\OLMathNumeral{2}}} \fCenter
    \maeh{{\OLId{greek-variable}{Delta}}_{\OLMathNumeral{1}}, !A({\OLId{latin-ordinary}{c}})}{{\OLId{greek-variable}{Delta}}_{\OLMathNumeral{2}}}$
    \RightLabel{\RightR{\lforall}}
    \UnaryInf$\maeh{{\OLId{greek-variable}{Gamma}}_{\OLMathNumeral{1}}}{{\OLId{greek-variable}{Gamma}}_{\OLMathNumeral{2}}} \fCenter
      \maeh{{\OLId{greek-variable}{Delta}}_{\OLMathNumeral{1}}, \lforall[{\OLId{latin-ordinary}{x}}][!A({\OLId{latin-ordinary}{x}})]}{{\OLId{greek-variable}{Delta}}_{\OLMathNumeral{2}}}$
    \DisplayProof
    \]
    أو
    \[
    \AxiomC{}
    \RightLabel{$\pi_{\OLMathNumeral{1}}$}
    \Deduce$\maeh{{\OLId{greek-variable}{Gamma}}_{\OLMathNumeral{1}}}{{\OLId{greek-variable}{Gamma}}_{\OLMathNumeral{2}}} \fCenter
      \maeh{{\OLId{greek-variable}{Delta}}_{\OLMathNumeral{1}}}{{\OLId{greek-variable}{Delta}}_{\OLMathNumeral{2}}, !A({\OLId{latin-ordinary}{c}})}$
    \RightLabel{\RightR{\lforall}}
    \UnaryInf$\maeh{{\OLId{greek-variable}{Gamma}}_{\OLMathNumeral{1}}}{{\OLId{greek-variable}{Gamma}}_{\OLMathNumeral{2}}} \fCenter
      \maeh{{\OLId{greek-variable}{Delta}}_{\OLMathNumeral{1}}}{{\OLId{greek-variable}{Delta}}_{\OLMathNumeral{2}}, \lforall[{\OLId{latin-ordinary}{x}}][!A({\OLId{latin-ordinary}{x}})]}$
    \DisplayProof
    \]
    في الوضع الأول يعطينا فرض الاستقراء~!!{proof}s لـ
    ${\OLId{greek-variable}{Gamma}}_{\OLMathNumeral{1}} \Sequent {\OLId{greek-variable}{Delta}}_{\OLMathNumeral{1}}, !A({\OLId{latin-ordinary}{c}}), !C_{\OLMathNumeral{1}}$ و$!C_{\OLMathNumeral{1}}, {\OLId{greek-variable}{Gamma}}_{\OLMathNumeral{2}}
    \Sequent {\OLId{greek-variable}{Delta}}_{\OLMathNumeral{2}}$. ويمكن تطبيق~$\RightR{\lforall}$ على الأولى
    فنحصل على~${\OLId{greek-variable}{Gamma}}_{\OLMathNumeral{1}} \Sequent {\OLId{greek-variable}{Delta}}_{\OLMathNumeral{1}}, \lforall[{\OLId{latin-ordinary}{x}}][!A({\OLId{latin-ordinary}{x}})], !C_{\OLMathNumeral{1}}$.
    (شرط المتغير المميز متحقق لأنه كان متحققًا أصلًا في قاعدة
    $\RightR{\lforall}$ في نهاية~$\pi$.) وعلى هذا تكون~$!C_{\OLMathNumeral{1}}$ صيغة
    استيفاء إذا تحقق الشرط~$\Lang L(!C_{\OLMathNumeral{1}}) \subseteq \Lang L({\OLId{greek-variable}{Gamma}}_{\OLMathNumeral{1}},
    {\OLId{greek-variable}{Delta}}_{\OLMathNumeral{1}}, \lforall[{\OLId{latin-ordinary}{x}}][!A({\OLId{latin-ordinary}{x}})]) \cap \Lang L({\OLId{greek-variable}{Gamma}}_{\OLMathNumeral{2}}, {\OLId{greek-variable}{Delta}}_{\OLMathNumeral{2}})$.
    ويعطينا فرض الاستقراء~$\Lang L(!C_{\OLMathNumeral{1}}) \subseteq \Lang L({\OLId{greek-variable}{Gamma}}_{\OLMathNumeral{1}},
    {\OLId{greek-variable}{Delta}}_{\OLMathNumeral{1}}, !A({\OLId{latin-ordinary}{c}})) \cap \Lang L({\OLId{greek-variable}{Gamma}}_{\OLMathNumeral{2}}, {\OLId{greek-variable}{Delta}}_{\OLMathNumeral{2}})$. فلذلك لا يبقى
    إلا القلق بشأن~${\OLId{latin-ordinary}{c}}$: هل يجوز أن تظهر في~$!C_{\OLMathNumeral{1}}$؟ لا يمكن. فلو ظهرت
    فيها لكان كل من~${\OLId{latin-ordinary}{c}} \in \Lang L({\OLId{greek-variable}{Gamma}}_{\OLMathNumeral{1}}, {\OLId{greek-variable}{Delta}}_{\OLMathNumeral{1}}, !A({\OLId{latin-ordinary}{c}}))$، وهو
    صحيح، و${\OLId{latin-ordinary}{c}} \in \Lang L({\OLId{greek-variable}{Gamma}}_{\OLMathNumeral{2}}, {\OLId{greek-variable}{Delta}}_{\OLMathNumeral{2}})$. غير أن~${\OLId{latin-ordinary}{c}}$ كانت عندئذ
    ستظهر في نتيجة استدلال~$\RightR{\lforall}$ في~$\pi$، خلافًا لشرط
    المتغير المميز.

    والحالة الأخرى مماثلة.

    \item ينتهي~$\pi$ بـ$\RightR{\lexists}$، وصيغته الرئيسة
    !!{formula}~$\lexists[{\OLId{latin-ordinary}{x}}][!A({\OLId{latin-ordinary}{x}})]$. ولدينا وضعان مرة أخرى. نتأمل
    الوضع الذي تنتمي فيه~$\lexists[{\OLId{latin-ordinary}{x}}][!A({\OLId{latin-ordinary}{x}})]$ إلى الجزء الأول، ونترك
    الآخر تمرينًا.

    في الوضع الأول ينتهي~$\pi$ بالشكل
    \[
    \AxiomC{}
    \RightLabel{$\pi_{\OLMathNumeral{1}}$}
    \Deduce$\maeh{{\OLId{greek-variable}{Gamma}}_{\OLMathNumeral{1}}}{{\OLId{greek-variable}{Gamma}}_{\OLMathNumeral{2}}} \fCenter
      \maeh{{\OLId{greek-variable}{Delta}}_{\OLMathNumeral{1}}, \lexists[{\OLId{latin-ordinary}{x}}][!A({\OLId{latin-ordinary}{x}})], !A({\OLId{latin-ordinary}{t}})}{{\OLId{greek-variable}{Delta}}_{\OLMathNumeral{2}}}$
    \RightLabel{\RightR{\lexists}}
    \UnaryInf$\maeh{{\OLId{greek-variable}{Gamma}}_{\OLMathNumeral{1}}}{{\OLId{greek-variable}{Gamma}}_{\OLMathNumeral{2}}} \fCenter
      \maeh{{\OLId{greek-variable}{Delta}}_{\OLMathNumeral{1}}, \lexists[{\OLId{latin-ordinary}{x}}][!A({\OLId{latin-ordinary}{x}})]}{{\OLId{greek-variable}{Delta}}_{\OLMathNumeral{2}}}$
    \DisplayProof
    \]
    ويعطينا فرض الاستقراء~!!{proof}s لـ${\OLId{greek-variable}{Gamma}}_{\OLMathNumeral{1}} \Sequent
    {\OLId{greek-variable}{Delta}}_{\OLMathNumeral{1}}, \lexists[{\OLId{latin-ordinary}{x}}][!A({\OLId{latin-ordinary}{x}})], !A({\OLId{latin-ordinary}{t}}), !C_{\OLMathNumeral{1}}$ و$!C_{\OLMathNumeral{1}}, {\OLId{greek-variable}{Gamma}}_{\OLMathNumeral{2}}
    \Sequent {\OLId{greek-variable}{Delta}}_{\OLMathNumeral{2}}$. ويمكن تطبيق~$\RightR{\lexists}$ على الأولى
    فنحصل على~${\OLId{greek-variable}{Gamma}}_{\OLMathNumeral{1}} \Sequent {\OLId{greek-variable}{Delta}}_{\OLMathNumeral{1}}, \lexists[{\OLId{latin-ordinary}{x}}][!A({\OLId{latin-ordinary}{x}})], !C_{\OLMathNumeral{1}}$.
    وعلى هذا تكون~$!C_{\OLMathNumeral{1}}$ صيغة استيفاء مرة أخرى إذا تحقق الشرط
    $\Lang L(!C_{\OLMathNumeral{1}}) \subseteq \Lang L({\OLId{greek-variable}{Gamma}}_{\OLMathNumeral{1}}, {\OLId{greek-variable}{Delta}}_{\OLMathNumeral{1}},
    \lexists[{\OLId{latin-ordinary}{x}}][!A({\OLId{latin-ordinary}{x}})]) \cap \Lang L({\OLId{greek-variable}{Gamma}}_{\OLMathNumeral{2}}, {\OLId{greek-variable}{Delta}}_{\OLMathNumeral{2}})$. أما فرض
    الاستقراء فيعطينا~$\Lang L(!C_{\OLMathNumeral{1}}) \subseteq \Lang L({\OLId{greek-variable}{Gamma}}_{\OLMathNumeral{1}},
    {\OLId{greek-variable}{Delta}}_{\OLMathNumeral{1}}, \lexists[{\OLId{latin-ordinary}{x}}][!A({\OLId{latin-ordinary}{x}})], !A({\OLId{latin-ordinary}{t}})) \cap \Lang L({\OLId{greek-variable}{Gamma}}_{\OLMathNumeral{2}},
    {\OLId{greek-variable}{Delta}}_{\OLMathNumeral{2}})$.

    ولأننا افترضنا عدم وجود~!!{function}s، فلا يجوز أن يكون الحد~${\OLId{latin-ordinary}{t}}$
    إلا متغيرًا أو~!!a{constant}. إذا كان متغيرًا فلا يضيف رمزًا إلى
    اللغة، فنأخذ~$!C_{\OLMathNumeral{1}}$ صيغة للاستيفاء. وإذا كان ثابتًا، ولنكتب
    ${\OLId{latin-ordinary}{t}} \ident {\OLId{latin-ordinary}{b}}$، فلا مشكلة كذلك إذا كان~${\OLId{latin-ordinary}{b}} \notin \Lang L(!C_{\OLMathNumeral{1}})$ أو
    كان~${\OLId{latin-ordinary}{b}} \in \Lang L({\OLId{greek-variable}{Gamma}}_{\OLMathNumeral{1}}, {\OLId{greek-variable}{Delta}}_{\OLMathNumeral{1}},
    \lexists[{\OLId{latin-ordinary}{x}}][!A({\OLId{latin-ordinary}{x}})]) \cap \Lang L({\OLId{greek-variable}{Gamma}}_{\OLMathNumeral{2}}, {\OLId{greek-variable}{Delta}}_{\OLMathNumeral{2}})$؛ فنأخذ
    $!C_{\OLMathNumeral{1}}$ صيغة للاستيفاء. لا نحتاج إلى التجريد إلا إذا كان~${\OLId{latin-ordinary}{b}} \in
    \Lang L(!C_{\OLMathNumeral{1}})$ ولم يكن~${\OLId{latin-ordinary}{b}}$ مشتركًا بين لغتي الجزأين. عندئذ نستبدل
    كل ظهور لـ${\OLId{latin-ordinary}{b}}$ في~$!C_{\OLMathNumeral{1}}$ بـ!!a{variable}~${\OLId{latin-ordinary}{y}}$، فنكتب
    $!C_{\OLMathNumeral{1}} \ident \Subst{!C_{\OLMathNumeral{1}}'}{{\OLId{latin-ordinary}{b}}}{{\OLId{latin-ordinary}{y}}}$، حيث لا تحتوي~$!C_{\OLMathNumeral{1}}'$ على~${\OLId{latin-ordinary}{b}}$.
    وإضافة إلى ذلك، إما أن يكون~${\OLId{latin-ordinary}{b}} \notin \Lang L({\OLId{greek-variable}{Gamma}}_{\OLMathNumeral{1}}, {\OLId{greek-variable}{Delta}}_{\OLMathNumeral{1}},
    \lexists[{\OLId{latin-ordinary}{x}}][!A({\OLId{latin-ordinary}{x}})])$ أو~${\OLId{latin-ordinary}{b}} \notin \Lang L({\OLId{greek-variable}{Gamma}}_{\OLMathNumeral{2}}, {\OLId{greek-variable}{Delta}}_{\OLMathNumeral{2}})$.
    وفي الحالتين نأخذ، على الترتيب،~$!C \ident
    \lforall[{\OLId{latin-ordinary}{y}}][!C_{\OLMathNumeral{1}}'({\OLId{latin-ordinary}{y}})]$ أو~$!C \ident \lexists[{\OLId{latin-ordinary}{y}}][!C_{\OLMathNumeral{1}}'({\OLId{latin-ordinary}{y}})]$.
    ففي الحالة الأولى لدينا:
    \[
    \AxiomC{}
    \Deduce${\OLId{greek-variable}{Gamma}}_{\OLMathNumeral{1}} \fCenter
      {\OLId{greek-variable}{Delta}}_{\OLMathNumeral{1}}, \lexists[{\OLId{latin-ordinary}{x}}][!A({\OLId{latin-ordinary}{x}})], !A({\OLId{latin-ordinary}{b}}), !C_{\OLMathNumeral{1}}'({\OLId{latin-ordinary}{b}})$
    \RightLabel{\RightR{\lexists}}
    \UnaryInf${\OLId{greek-variable}{Gamma}}_{\OLMathNumeral{1}} \fCenter {\OLId{greek-variable}{Delta}}_{\OLMathNumeral{1}}, \lexists[{\OLId{latin-ordinary}{x}}][!A({\OLId{latin-ordinary}{x}})], !C_{\OLMathNumeral{1}}'({\OLId{latin-ordinary}{b}})$
    \RightLabel{\RightR{\lforall}}
    \UnaryInf${\OLId{greek-variable}{Gamma}}_{\OLMathNumeral{1}} \fCenter
      {\OLId{greek-variable}{Delta}}_{\OLMathNumeral{1}}, \lexists[{\OLId{latin-ordinary}{x}}][!A({\OLId{latin-ordinary}{x}})], \lforall[{\OLId{latin-ordinary}{y}}][!C_{\OLMathNumeral{1}}'({\OLId{latin-ordinary}{y}})]$
    \DisplayProof
    \]
    و
    \[
    \AxiomC{}
    \Deduce$!C_{\OLMathNumeral{1}}'({\OLId{latin-ordinary}{b}}), \lforall[{\OLId{latin-ordinary}{y}}][!C_{\OLMathNumeral{1}}'({\OLId{latin-ordinary}{y}})], {\OLId{greek-variable}{Gamma}}_{\OLMathNumeral{2}} \fCenter {\OLId{greek-variable}{Delta}}_{\OLMathNumeral{2}}$
    \RightLabel{\LeftR{\lforall}}
    \UnaryInf$\lforall[{\OLId{latin-ordinary}{y}}][!C_{\OLMathNumeral{1}}'({\OLId{latin-ordinary}{y}})], {\OLId{greek-variable}{Gamma}}_{\OLMathNumeral{2}} \fCenter {\OLId{greek-variable}{Delta}}_{\OLMathNumeral{2}}$
    \DisplayProof
    \]
    وتأتي~!!{proof}s المتتاليتين العلويتين من فرض الاستقراء (وفي
    الثانية نطبق مقبولية الإضعاف مع~$\lforall[{\OLId{latin-ordinary}{y}}][!C_{\OLMathNumeral{1}}'({\OLId{latin-ordinary}{y}})]$ في
    المقدّم). ويتحقق شرط المتغير المميز في~$\RightR{\lforall}$ في
    !!{proof} الأول لأن~${\OLId{latin-ordinary}{b}} \notin \Lang L({\OLId{greek-variable}{Gamma}}_{\OLMathNumeral{1}}, {\OLId{greek-variable}{Delta}}_{\OLMathNumeral{1}},
    \lexists[{\OLId{latin-ordinary}{x}}][!A({\OLId{latin-ordinary}{x}})])$، ولأن~$!C_{\OLMathNumeral{1}}'$ عُرّفت بحيث لا تحتوي~${\OLId{latin-ordinary}{b}}$.

    وفي الحالة الثانية لدينا، من فرض الاستقراء مع مقبولية الإضعاف بـ
    $\lexists[{\OLId{latin-ordinary}{y}}][!C_{\OLMathNumeral{1}}'({\OLId{latin-ordinary}{y}})]$:
    \[
    \AxiomC{}
    \Deduce${\OLId{greek-variable}{Gamma}}_{\OLMathNumeral{1}} \fCenter
      {\OLId{greek-variable}{Delta}}_{\OLMathNumeral{1}}, \lexists[{\OLId{latin-ordinary}{x}}][!A({\OLId{latin-ordinary}{x}})], !A({\OLId{latin-ordinary}{b}}), \lexists[{\OLId{latin-ordinary}{y}}][!C_{\OLMathNumeral{1}}'({\OLId{latin-ordinary}{y}})], !C_{\OLMathNumeral{1}}'({\OLId{latin-ordinary}{b}})$
    \RightLabel{\RightR{\lexists}}
    \UnaryInf${\OLId{greek-variable}{Gamma}}_{\OLMathNumeral{1}} \fCenter
      {\OLId{greek-variable}{Delta}}_{\OLMathNumeral{1}}, \lexists[{\OLId{latin-ordinary}{x}}][!A({\OLId{latin-ordinary}{x}})], \lexists[{\OLId{latin-ordinary}{y}}][!C_{\OLMathNumeral{1}}'({\OLId{latin-ordinary}{y}})], !C_{\OLMathNumeral{1}}'({\OLId{latin-ordinary}{b}})$
    \RightLabel{\RightR{\lexists}}
    \UnaryInf${\OLId{greek-variable}{Gamma}}_{\OLMathNumeral{1}} \fCenter
      {\OLId{greek-variable}{Delta}}_{\OLMathNumeral{1}}, \lexists[{\OLId{latin-ordinary}{x}}][!A({\OLId{latin-ordinary}{x}})], \lexists[{\OLId{latin-ordinary}{y}}][!C_{\OLMathNumeral{1}}'({\OLId{latin-ordinary}{y}})]$
    \DisplayProof
    \]
    و
    \[
    \AxiomC{}
    \Deduce$!C_{\OLMathNumeral{1}}'({\OLId{latin-ordinary}{b}}), {\OLId{greek-variable}{Gamma}}_{\OLMathNumeral{2}} \fCenter {\OLId{greek-variable}{Delta}}_{\OLMathNumeral{2}}$
    \RightLabel{\LeftR{\lexists}}
    \UnaryInf$\lexists[{\OLId{latin-ordinary}{y}}][!C_{\OLMathNumeral{1}}'({\OLId{latin-ordinary}{y}})], {\OLId{greek-variable}{Gamma}}_{\OLMathNumeral{2}} \fCenter {\OLId{greek-variable}{Delta}}_{\OLMathNumeral{2}}$
    \DisplayProof
    \]
    ويتحقق شرط المتغير المميز في~$\LeftR{\lexists}$ في~!!{proof}
    الثاني لأن~${\OLId{latin-ordinary}{b}} \notin \Lang L({\OLId{greek-variable}{Gamma}}_{\OLMathNumeral{2}}, {\OLId{greek-variable}{Delta}}_{\OLMathNumeral{2}})$، ولأن~$!C_{\OLMathNumeral{1}}'$
    عُرّفت بحيث لا تحتوي~${\OLId{latin-ordinary}{b}}$.
  \end{enumerate}
  والحالات الأخرى مماثلة، ونتركها تمارين.
\end{proof}

\begin{prob}
لنفرض أن لدينا الرابط~$\lnand$ ذا القاعدتين
\[
\Axiom${\OLId{greek-variable}{Gamma}} \fCenter {\OLId{greek-variable}{Delta}}, !A$
\Axiom${\OLId{greek-variable}{Gamma}} \fCenter {\OLId{greek-variable}{Delta}}, !B$
\RightLabel{\LeftR{\lnand}}
\BinaryInf$!A \lnand !B, {\OLId{greek-variable}{Gamma}} \fCenter {\OLId{greek-variable}{Delta}}$
\DisplayProof
\quad
\Axiom$!A, !B, {\OLId{greek-variable}{Gamma}} \fCenter {\OLId{greek-variable}{Delta}}$
\RightLabel{\RightR{\lnand}}
\UnaryInf${\OLId{greek-variable}{Gamma}} \fCenter {\OLId{greek-variable}{Delta}}, !A \lnand !B$
\DisplayProof
\]
بيّن أن لمّة ماهارا تظل صادقة بمعالجة الحالات في برهان
\olref{lem:maehara} التي ينتهي فيها~$\pi$ بإحدى هاتين القاعدتين.
\end{prob}

بعد إثبات لمّة ماهارا، نستطيع استخدامها لبرهنة مبرهنة كريغ للاستيفاء.

\begin{proof}[برهان \olref{thm:interpolation}]
  لنفرض أن~$!A \Sequent !B$~!!{provable}. طبّق~\olref{lem:maehara}
  على~$\maeh{!A}{\ } \Sequent \maeh{\ }{!B}$ لتحصل على صيغة
  الاستيفاء~$!C$. وعندئذ~$\Proves !A \Sequent !C$ و$\Proves !C
  \Sequent !B$.
\end{proof}

لنفرض أن~${\OLId{greek-variable}{Gamma}}$ نظرية (مجموعة~!!{sentence}s) في
!!a{language}~$\Lang L$ تحتوي~!!{predicate}~${\OLId{latin-looped}{R}}$ ذا~${\OLId{latin-ordinary}{n}}$ مواضع. نقول
إن~${\OLId{greek-variable}{Gamma}}$ \emph{تعرّف}~${\OLId{latin-looped}{R}}$ \emph{ضمنيًا} إذا وفقط إذا
\begin{align*}
{\OLId{greek-variable}{Gamma}}, {\OLId{greek-variable}{Gamma}}'
& \Entails \lforall[{\OLId{latin-ordinary}{x}}_{\OLMathNumeral{1}}][\dots\lforall[{\OLId{latin-ordinary}{x}}_{\OLId{latin-ordinary}{n}}][({\OLId{latin-looped}{R}}({\OLId{latin-ordinary}{x}}_{\OLMathNumeral{1}},\dots,{\OLId{latin-ordinary}{x}}_{\OLId{latin-ordinary}{n}}) \liff
{\OLId{latin-looped}{R}}'({\OLId{latin-ordinary}{x}}_{\OLMathNumeral{1}}, \dots, {\OLId{latin-ordinary}{x}}_{\OLId{latin-ordinary}{n}}))]],
\intertext{حيث~${\OLId{greek-variable}{Gamma}}'$ هي~${\OLId{greek-variable}{Gamma}}$ بعد استبدال كل ظهور لـ${\OLId{latin-looped}{R}}$ بـ
${\OLId{latin-looped}{R}}' \notin \Lang L$. وتكون~${\OLId{greek-variable}{Gamma}}$ قد \emph{عرّفت}~${\OLId{latin-looped}{R}}$
\emph{صراحة} إذا وجدت~$!A({\OLId{latin-ordinary}{x}}_{\OLMathNumeral{1}}, \dots, {\OLId{latin-ordinary}{x}}_{\OLId{latin-ordinary}{n}})$ لا تحتوي~${\OLId{latin-looped}{R}}$ بحيث}
{\OLId{greek-variable}{Gamma}} &\Entails
\lforall[{\OLId{latin-ordinary}{x}}_{\OLMathNumeral{1}}][\dots\lforall[{\OLId{latin-ordinary}{x}}_{\OLId{latin-ordinary}{n}}][({\OLId{latin-looped}{R}}({\OLId{latin-ordinary}{x}}_{\OLMathNumeral{1}},\dots,{\OLId{latin-ordinary}{x}}_{\OLId{latin-ordinary}{n}}) \liff !A({\OLId{latin-ordinary}{x}}_{\OLMathNumeral{1}},
\dots, {\OLId{latin-ordinary}{x}}_{\OLId{latin-ordinary}{n}}))]].
\end{align*}
ومن الواضح أن تعريف~${\OLId{greek-variable}{Gamma}}$ لـ${\OLId{latin-looped}{R}}$ صراحة يستلزم تعريفها له ضمنيًا.
أما العكس فليس واضحًا، لكنه يصدق مع ذلك:

\begin{lem}[مبرهنة بيث في قابلية التعريف]
  إذا عرّفت~${\OLId{greek-variable}{Gamma}}$ العلاقة~${\OLId{latin-looped}{R}}$ ضمنيًا، عرّفتها صراحة كذلك.
\end{lem}

\begin{proof}
  لنفرض أن~${\OLId{greek-variable}{Gamma}}$ تعرّف~${\OLId{latin-looped}{R}}$ ضمنيًا. لتكن~${\OLId{latin-looped}{R}}'$ و${\OLId{greek-variable}{Gamma}}'$ كما سبق،
  ولتكن~$\Lang L' = \Lang L({\OLId{greek-variable}{Gamma}}') = \Lang L \setminus \{{\OLId{latin-looped}{R}}\}
  \cup \{{\OLId{latin-looped}{R}}'\}$. وبالفرض لدينا
  \begin{align*}
    {\OLId{greek-variable}{Gamma}}, {\OLId{greek-variable}{Gamma}}' &
    \Sequent \lforall[{\OLId{latin-ordinary}{x}}_{\OLMathNumeral{1}}][\dots\lforall[{\OLId{latin-ordinary}{x}}_{\OLId{latin-ordinary}{n}}][({\OLId{latin-looped}{R}}({\OLId{latin-ordinary}{x}}_{\OLMathNumeral{1}},\dots,{\OLId{latin-ordinary}{x}}_{\OLId{latin-ordinary}{n}})
    \liff {\OLId{latin-looped}{R}}'({\OLId{latin-ordinary}{x}}_{\OLMathNumeral{1}}, \dots, {\OLId{latin-ordinary}{x}}_{\OLId{latin-ordinary}{n}}))]].
  \intertext{لتكن~${\OLId{latin-ordinary}{c}}_{\OLMathNumeral{1}}$, \dots,~${\OLId{latin-ordinary}{c}}_{\OLId{latin-ordinary}{n}}$~!!{constant}s لا تظهر في
  ${\OLId{greek-variable}{Gamma}}$. وتعطينا لمّة القابلية للعكس}
    {\OLId{greek-variable}{Gamma}}, {\OLId{greek-variable}{Gamma}}', {\OLId{latin-looped}{R}}({\OLId{latin-ordinary}{c}}_{\OLMathNumeral{1}},\dots,{\OLId{latin-ordinary}{c}}_{\OLId{latin-ordinary}{n}}) & \Sequent  {\OLId{latin-looped}{R}}'({\OLId{latin-ordinary}{c}}_{\OLMathNumeral{1}}, \dots, {\OLId{latin-ordinary}{c}}_{\OLId{latin-ordinary}{n}}).
  \intertext{عرّف}
    \Lang L_{\OLMathNumeral{1}} &= \Lang L({\OLId{greek-variable}{Gamma}}, {\OLId{latin-looped}{R}}({\OLId{latin-ordinary}{c}}_{\OLMathNumeral{1}},\dots,{\OLId{latin-ordinary}{c}}_{\OLId{latin-ordinary}{n}})), &
    \Lang L_{\OLMathNumeral{2}} &= \Lang L({\OLId{greek-variable}{Gamma}}',{\OLId{latin-looped}{R}}'({\OLId{latin-ordinary}{c}}_{\OLMathNumeral{1}},\dots,{\OLId{latin-ordinary}{c}}_{\OLId{latin-ordinary}{n}})).
  \intertext{وطبّق لمّة ماهارا على}
    \maeh{{\OLId{greek-variable}{Gamma}}, {\OLId{latin-looped}{R}}({\OLId{latin-ordinary}{c}}_{\OLMathNumeral{1}},\dots,{\OLId{latin-ordinary}{c}}_{\OLId{latin-ordinary}{n}})}{{\OLId{greek-variable}{Gamma}}'} & \Sequent
      \maeh{\ }{{\OLId{latin-looped}{R}}'({\OLId{latin-ordinary}{c}}_{\OLMathNumeral{1}}, \dots, {\OLId{latin-ordinary}{c}}_{\OLId{latin-ordinary}{n}})}.
  \intertext{فنحصل على صيغة استيفاء~$!C({\OLId{latin-ordinary}{c}}_{\OLMathNumeral{1}},\dots,{\OLId{latin-ordinary}{c}}_{\OLId{latin-ordinary}{n}})$ بحيث}
    {\OLId{greek-variable}{Gamma}}, {\OLId{latin-looped}{R}}({\OLId{latin-ordinary}{c}}_{\OLMathNumeral{1}},\dots,{\OLId{latin-ordinary}{c}}_{\OLId{latin-ordinary}{n}}) & \Sequent !C({\OLId{latin-ordinary}{c}}_{\OLMathNumeral{1}}, \dots, {\OLId{latin-ordinary}{c}}_{\OLId{latin-ordinary}{n}}) \\
    !C({\OLId{latin-ordinary}{c}}_{\OLMathNumeral{1}}, \dots, {\OLId{latin-ordinary}{c}}_{\OLId{latin-ordinary}{n}}), {\OLId{greek-variable}{Gamma}}' & \Sequent {\OLId{latin-looped}{R}}'({\OLId{latin-ordinary}{c}}_{\OLMathNumeral{1}}, \dots, {\OLId{latin-ordinary}{c}}_{\OLId{latin-ordinary}{n}}).
  \intertext{ولهاتين المتتاليتين~!!{proof}s. وتعطينا لمّة ماهارا}
    \Lang L(!C) &\subseteq \Lang L_{\OLMathNumeral{1}} \cap \Lang L_{\OLMathNumeral{2}} \\
      &= (\Lang L({\OLId{greek-variable}{Gamma}}) \setminus \{{\OLId{latin-looped}{R}}\})
         \cup \{{\OLId{latin-ordinary}{c}}_{\OLMathNumeral{1}},\dots,{\OLId{latin-ordinary}{c}}_{\OLId{latin-ordinary}{n}}\}.
  \intertext{لتكن~$!A({\OLId{latin-ordinary}{x}}_{\OLMathNumeral{1}},\dots,{\OLId{latin-ordinary}{x}}_{\OLId{latin-ordinary}{n}})$ ناتجة من~$!C({\OLId{latin-ordinary}{c}}_{\OLMathNumeral{1}},\dots,{\OLId{latin-ordinary}{c}}_{\OLId{latin-ordinary}{n}})$
  باستبدال الثابت الجديد~${\OLId{latin-ordinary}{c}}_{\OLId{latin-ordinary}{i}}$ بالمتغير~${\OLId{latin-ordinary}{x}}_{\OLId{latin-ordinary}{i}}$ لكل~${\OLId{latin-ordinary}{i}}$. عندئذ لا تحتوي
  $!A$ الرمز~${\OLId{latin-looped}{R}}$. واستبدل~${\OLId{latin-looped}{R}}'$ في كل موضع بـ${\OLId{latin-looped}{R}}$ في~!!{proof} المتتالية
  الثانية لتحصل على~!!a{proof} لـ}
    !A({\OLId{latin-ordinary}{c}}_{\OLMathNumeral{1}}, \dots, {\OLId{latin-ordinary}{c}}_{\OLId{latin-ordinary}{n}}), {\OLId{greek-variable}{Gamma}} & \Sequent {\OLId{latin-looped}{R}}({\OLId{latin-ordinary}{c}}_{\OLMathNumeral{1}}, \dots, {\OLId{latin-ordinary}{c}}_{\OLId{latin-ordinary}{n}}).
  \end{align*}
  والمتتالية الأولى هي
  ${\OLId{greek-variable}{Gamma}}, {\OLId{latin-looped}{R}}({\OLId{latin-ordinary}{c}}_{\OLMathNumeral{1}},\dots,{\OLId{latin-ordinary}{c}}_{\OLId{latin-ordinary}{n}}) \Sequent !A({\OLId{latin-ordinary}{c}}_{\OLMathNumeral{1}},\dots,{\OLId{latin-ordinary}{c}}_{\OLId{latin-ordinary}{n}})$.
  وبالجمع القضوي بين المتتاليتين ينتج لنا
  \[
    {\OLId{greek-variable}{Gamma}} \Sequent {\OLId{latin-looped}{R}}({\OLId{latin-ordinary}{c}}_{\OLMathNumeral{1}},\dots,{\OLId{latin-ordinary}{c}}_{\OLId{latin-ordinary}{n}}) \liff !A({\OLId{latin-ordinary}{c}}_{\OLMathNumeral{1}},\dots,{\OLId{latin-ordinary}{c}}_{\OLId{latin-ordinary}{n}}).
  \]
  ثم نطبق~$\RightR{\lforall}$ باستعمال الثوابت الجديدة تباعًا فنحصل على
  !!a{proof} لـ
  \[
    {\OLId{greek-variable}{Gamma}} \Sequent \lforall[{\OLId{latin-ordinary}{x}}_{\OLMathNumeral{1}}][\dots\lforall[{\OLId{latin-ordinary}{x}}_{\OLId{latin-ordinary}{n}}][({\OLId{latin-looped}{R}}({\OLId{latin-ordinary}{x}}_{\OLMathNumeral{1}},\dots,{\OLId{latin-ordinary}{x}}_{\OLId{latin-ordinary}{n}}) \liff !A({\OLId{latin-ordinary}{x}}_{\OLMathNumeral{1}}, \dots, {\OLId{latin-ordinary}{x}}_{\OLId{latin-ordinary}{n}}))]].
  \]
  وهذا يبيّن أن~$!A({\OLId{latin-ordinary}{x}}_{\OLMathNumeral{1}}, \dots, {\OLId{latin-ordinary}{x}}_{\OLId{latin-ordinary}{n}})$ تعرّف~${\OLId{latin-looped}{R}}$ صراحة في~${\OLId{greek-variable}{Gamma}}$.
\end{proof}

وثمة نتيجة ثالثة تكافئ الاستيفاء، ويمكن مثلها برهنتها باستعمال لمّة
ماهارا، وهي: إذا كانت~${\OLId{greek-variable}{Gamma}}_{\OLMathNumeral{1}}$ و${\OLId{greek-variable}{Gamma}}_{\OLMathNumeral{2}}$ نظريتين متسقتين تحتويان
نظرية تامة~${\OLId{greek-variable}{Gamma}}$ في اللغة~$\Lang L({\OLId{greek-variable}{Gamma}}) \subseteq
\Lang L({\OLId{greek-variable}{Gamma}}_{\OLMathNumeral{1}}) \cap \Lang L({\OLId{greek-variable}{Gamma}}_{\OLMathNumeral{2}})$، كانت
${\OLId{greek-variable}{Gamma}}_{\OLMathNumeral{1}} \cup {\OLId{greek-variable}{Gamma}}_{\OLMathNumeral{2}}$ متسقة.

تذكّر أن النظرية التامة هي التي يكون فيها، لكل~!!{sentence}~$!A$ في
لغتها، إما~${\OLId{greek-variable}{Gamma}} \Proves !A$ أو~${\OLId{greek-variable}{Gamma}} \Proves \lnot !A$. ولأن
${\OLId{greek-variable}{Gamma}}_{\OLMathNumeral{1}}$ و${\OLId{greek-variable}{Gamma}}_{\OLMathNumeral{2}}$ امتدادان متسقان لـ${\OLId{greek-variable}{Gamma}}$، فلا بد أن تكون
${\OLId{greek-variable}{Gamma}}$ متسقة. ومن ثم، إذا كانت~$!A$~!!a{sentence} في اللغة
$\Lang L({\OLId{greek-variable}{Gamma}}_{\OLMathNumeral{1}}) \cap \Lang L({\OLId{greek-variable}{Gamma}}_{\OLMathNumeral{2}})$، امتنع أن يكون معًا
${\OLId{greek-variable}{Gamma}}_{\OLMathNumeral{1}} \Entails !A$ و${\OLId{greek-variable}{Gamma}}_{\OLMathNumeral{2}} \Entails \lnot !A$. ولرؤية ذلك،
افترض وجود~$!A$ كهذه. إذا كان~${\OLId{greek-variable}{Gamma}}_{\OLMathNumeral{1}} \Entails !A$، لزم أن يكون
${\OLId{greek-variable}{Gamma}} \Entails !A$؛ وإلا استلزم تمام~${\OLId{greek-variable}{Gamma}}$ أن
${\OLId{greek-variable}{Gamma}} \Entails \lnot !A$، وبما أن~${\OLId{greek-variable}{Gamma}} \subseteq {\OLId{greek-variable}{Gamma}}_{\OLMathNumeral{1}}$،
كان~${\OLId{greek-variable}{Gamma}}_{\OLMathNumeral{1}} \Entails \lnot !A$، فتغدو~${\OLId{greek-variable}{Gamma}}_{\OLMathNumeral{1}}$ غير متسقة. إذن
${\OLId{greek-variable}{Gamma}} \Entails !A$، وبما أن~${\OLId{greek-variable}{Gamma}} \subseteq {\OLId{greek-variable}{Gamma}}_{\OLMathNumeral{2}}$، يكون
${\OLId{greek-variable}{Gamma}}_{\OLMathNumeral{2}} \Entails !A$. ولذلك~${\OLId{greek-variable}{Gamma}}_{\OLMathNumeral{2}} \Entails/ \lnot !A$، وإلا
كانت~${\OLId{greek-variable}{Gamma}}_{\OLMathNumeral{2}}$ غير متسقة.

وعدم وجود~!!a{sentence}~$!A$ في~!!{language} المشتركة
$\Lang L({\OLId{greek-variable}{Gamma}}_{\OLMathNumeral{1}}) \cap \Lang L({\OLId{greek-variable}{Gamma}}_{\OLMathNumeral{2}})$ بحيث
${\OLId{greek-variable}{Gamma}}_{\OLMathNumeral{1}} \Entails !A$ و${\OLId{greek-variable}{Gamma}}_{\OLMathNumeral{2}} \Entails \lnot !A$ هو كل ما يلزم
للبرهان، الذي نقدمه هنا برهانيًا باستعمال لمّة ماهارا.

\begin{thm}[مبرهنة روبنسون في الاتساق المشترك]
لنفرض أن~${\OLId{greek-variable}{Gamma}}_{\OLMathNumeral{1}}$ و${\OLId{greek-variable}{Gamma}}_{\OLMathNumeral{2}}$ نظريتان متسقتان، وأنه لا توجد~$!A$ في
!!{language}~$\Lang L({\OLId{greek-variable}{Gamma}}_{\OLMathNumeral{1}}) \cap \Lang L({\OLId{greek-variable}{Gamma}}_{\OLMathNumeral{2}})$ بحيث يكون
${\OLId{greek-variable}{Gamma}}_{\OLMathNumeral{1}} \Entails !A$ و${\OLId{greek-variable}{Gamma}}_{\OLMathNumeral{2}} \Entails \lnot !A$. عندئذ تكون
${\OLId{greek-variable}{Gamma}}_{\OLMathNumeral{1}} \cup {\OLId{greek-variable}{Gamma}}_{\OLMathNumeral{2}}$ متسقة.
\end{thm}

\begin{proof}
  افترض، خلافًا للمطلوب، أن~${\OLId{greek-variable}{Gamma}}_{\OLMathNumeral{1}} \cup {\OLId{greek-variable}{Gamma}}_{\OLMathNumeral{2}}$ غير متسقة. عندئذ
  للمتتالية~${\OLId{greek-variable}{Gamma}}_{\OLMathNumeral{1}}, {\OLId{greek-variable}{Gamma}}_{\OLMathNumeral{2}} \Sequent \ $~!!a{proof}. طبّق لمّة
  ماهارا على~$\maeh{{\OLId{greek-variable}{Gamma}}_{\OLMathNumeral{1}}}{{\OLId{greek-variable}{Gamma}}_{\OLMathNumeral{2}}} \Sequent \maeh{\ }{\ }$
  لتحصل على~!!a{sentence}~$!A$ في~!!{language}~$\Lang L({\OLId{greek-variable}{Gamma}}_{\OLMathNumeral{1}})
  \cap \Lang L({\OLId{greek-variable}{Gamma}}_{\OLMathNumeral{2}})$ بحيث~$\Proves {\OLId{greek-variable}{Gamma}}_{\OLMathNumeral{1}} \Sequent !A$ و
  $\Proves !A, {\OLId{greek-variable}{Gamma}}_{\OLMathNumeral{2}} \Sequent \quad$. ومن الأخيرة ينتج لنا
  $\Proves {\OLId{greek-variable}{Gamma}}_{\OLMathNumeral{2}} \Sequent \lnot !A$. لكننا افترضنا عدم وجود
  !!{sentence} كهذه.
\end{proof}

افترضنا ضمنًا فيما سبق أن النظريات~${\OLId{greek-variable}{Gamma}}$ و${\OLId{greek-variable}{Gamma}}_{\OLMathNumeral{1}}$ و${\OLId{greek-variable}{Gamma}}_{\OLMathNumeral{2}}$
منتهية. وبفضل التراص تصدق النتائج كذلك على النظريات اللانهائية.

\end{document}
